%   The template should run in “xeLatex”  format. Put "ICN-font.otf" into .../MikTex/fonts/***/.
\documentclass[10.5pt,compsoc]{ICN}
\setcounter{page}{1}
\graphicspath{{figures/}}
\usepackage{xeCJK}
\usepackage{titlesec}
\usepackage[utf8]{inputenc}
\usepackage{fontenc,xunicode,xltxtra}
\usepackage{fontspec}
\usepackage{graphicx}
\usepackage{footmisc}
\usepackage{subfigure}
\usepackage{url}
\usepackage{multirow}
\usepackage[noadjust]{cite}
\usepackage{amsmath,amsthm}
\usepackage{amssymb,amsfonts}
\usepackage{unicode-math}
\usepackage{booktabs}
\usepackage{color}
\usepackage{ccaption}
\usepackage{booktabs}
\usepackage{float}
\usepackage{fancyhdr}
\usepackage{caption}
\usepackage{xcolor,stfloats}
\usepackage{comment}
\usepackage{cuted}  %flushend,
\usepackage{captionhack}
\usepackage{epstopdf}

%\usepackage{titlesec}%定义标题宏命令

%===============================%
\headevenname{\size{-5}{\textbf{\lwbt \textit{{Intelligent and Converged Networks,}}  xxxxxxx}}  20xx, x(x): xxx-xxx}
\headoddname{{First author et al.:}\quad {\textbf{\emph{Click and Type Your Title Here ...}}}
}%

\renewcommand{\thefootnote}{\fnsymbol{footnote}}
\setcounter{footnote}{0}
\renewcommand\footnotelayout{\normalsize}

\newtheoremstyle{mystyle}{0pt}{0pt}{\normalfont}{1em}{\bf}{}{1em}{}
\theoremstyle{mystyle}
\renewcommand\figurename{Fig.~}
\renewcommand{\thesubfigure}{(\alph{subfigure})}
\newcommand{\upcite}[1]{\textsuperscript{\cite{#1}}}
\renewcommand{\labelenumi}{(\arabic{enumi})}
\newcommand{\tabincell}[2]{\begin{tabular}{@{}#1@{}}#2\end{tabular}}
\newcommand{\abc}{\color{white}\vrule width 2pt}


\newtheorem{assumption}{\textbf{Assumption}}
\newtheorem{definition}{\textbf{Definition}}
\newtheorem{lemma}{\textbf{Lemma}}
\newtheorem{theorem}{\textbf{Theorem}}
\newtheorem{proposition}{\textbf{Proposition}}
\newtheorem{corollary}{\textbf{Corollary}}
\newcommand{\nop}[1]{}


\addtolength{\abovecaptionskip}{-2mm}
\addtolength{\belowcaptionskip}{-2mm}

\makeatletter
\renewcommand{\@biblabel}[1]{[#1]\hfill}
\makeatother


\linespread{1.2}
\setmainfont[BoldFont={CALIBRIB.TTF},ItalicFont={CALIBRII.TTF},BoldItalicFont={CALIBRIZ.TTF}]{CALIBRI.TTF}%
\setmathfont{Times New Roman}
%\setCJKmainfont[BoldFont=KaiTi,ItalicFont=KaiTi]{Calibri}
\setCJKsansfont{CALIBRI.TTF}
\setmonofont{CALIBRI.TTF}
\setsansfont{CALIBRI.TTF}
\newcommand{\tnewroman}{\fontspec{Times New Roman}}
\newcommand{\lwbt}{\fontspec{bahnschrift.ttf}}%{Lato Black}}
\newcommand{\lwbtI}{\fontspec{bahnschrift.ttf}}%{Lato Black Italic}}


%\newcommand{\lwbt}{\fontspec{ICN-font.otf}}%{Lato Black}}
%\newcommand{\lwbtI}{\fontspec{ICN-font.otf}}%{Lato Black Italic}}


\XeTeXlinebreaklocale "zh"
\XeTeXlinebreakskip = 0pt plus 1pt minus 0.1pt


\setCJKfamilyfont{song}{SimSun}
\newcommand{\song}{\CJKfamily{song}}
\setCJKfamilyfont{xs}{SimSun}
\newcommand{\xs}{\CJKfamily{xs}}
\begin{document}
%\raggedbottom
\thispagestyle{empty}
\begin{strip}\size{3}
\noindent
\\{\lwbt{Template for Preparation of Manuscripts for}\\[-1mm]\lwbtI{Intelligent and Converged Networks}}
\vskip +6mm
\size{5}
\noindent
This template is to be used for preparing manuscripts for submission to \emph{Intelligent and Converged Network}. Use of this template will save time in the review and production processes and will expedite publication. Do not modify the template in any way (delete spaces, modify font size/line height, etc.).

\vspace{180mm}
\end{strip}
\clearpage
\hyphenpenalty=50000
\makeatletter
\newcommand\mysmall{\@setfontsize\mysmall{7}{9.5}}
%%%%%%%%%%%%%%%%%%%%%%%
\newenvironment{tablehere}
  {\def\@captype{table}}
  {}
\newenvironment{figurehere}
  {\def\@captype{figure}}
  {}
%%%%%%%%%%%%%%%%%%%%%%%
\thispagestyle{plain}%
\thispagestyle{empty}%


\let\temp\footnote
\renewcommand \footnote[1]{\temp{\size{-5}#1}}
{}
\vspace*{-45pt}
\hspace{-95mm}
\noindent\begin{tabular}{p{70mm}p{100mm}}
\hspace{-5mm} \lwbt {\size{5-}\textbf{\scalebox{0.95}[1.0]{\makebox[5cm][s]{Intelligent and Converged Networks}}}}&\lwbt
\hfill{{\size{5-}\textbf{\scalebox{1}[1.0]{\makebox[5cm][s]{{Volume\hspace{0.356em}x,\hspace{0.356em}Number\hspace{0.356em}x,\hspace{0.356em}
xxxxxxx  20xx}}}}}}\\
{\size{5-}
\textbf{\lwbt
\hspace{-5mm}
\scalebox{1}[1.0]{\makebox[5.6cm][s]{%
ISSN xxxx-xxxx
\quad {\color{white}{0x/xx\quad ppxxx--xx x}\hfill}}}}}&
\lwbt \hfill\size{5-}{\textbf{\scalebox{1}[1.0]{\makebox[4.6cm][s]{
~DOI: 10.26599/ICN.20xx.90200xx}}}}\end{tabular}

\begin{strip}

{\centering
{\size{3}\textbf{\lwbt Click and Type Your Title Here. The Font is Bahnschrift \\[0mm]
$\backslash$size\{3\}. The Title Must be Brief and Grammatically Correct}}
\vskip 9mm}
{\center {\size{5}
Type author names here. The font is Calibri ``$\backslash$size\{5\}". Please spell out first names and surnames.  Place an ``$*$" by the corresponding author(s). For example, First A. Author, Second B. Author$^*$, and Third C. Author
}\vskip 5mm}


\centering
\begin{tabular}{p{160mm}}
{\setlength{\baselineskip}{18.5pt}\selectfont
\noindent
\textbf{Abstract:} {Type your abstract text here. The font is Calibri $\backslash$size\{5\--\}. All manuscripts must be accompanied by an abstract which briefly summarizes the reasons for the work, the methods, the main results and conclusions of the study. The optimal length is one paragraph (100-250 words). Since the abstract is indexed in many da-tabases it should be self-contained without any undefined abbreviations or any cited references.}
\vskip 4mm
\noindent
\textbf{ Key words:} {Type keywords (3 to 8 keywords) here. The font is Calibri $\backslash$size\{5\--\}. For example,  gene regulation; ...}\par}
\end{tabular}
\vskip 6mm
\vskip -3mm
\size{6}\end{strip}
\thispagestyle{plain}%
\thispagestyle{empty}%
\makeatother
\pagestyle{tstheadings}
\begin{figure}[b]
\vskip -3mm
\begin{tabular}{p{44mm}}
\toprule\\
\end{tabular}
\vskip -5.5mm
\noindent
\setlength{\tabcolsep}{1pt}
\begin{tabular}{p{1.5mm}p{79.5mm}}
&Type the author's affiliation together with E-mail address here. The font is Calibri $\backslash$size\{6\}. For example,  \\
%{$\sf{\hskip 0.5em*}$}
$\bullet$& First A. Author and Second B. Author are with the Department, University, City and Postcode, Country. E-mail:\\
$\bullet$& Third C. Author is with the Department, University, City and Postcode, Country. E-mail:
 \\
$*$&To whom correspondence should be addressed. \\
         &Manuscript received: year-month-day;
          %revised: year-month-day;
          accepted: year-month-day
\end{tabular}
\vspace {0mm}
\end{figure}
\size{5}
\vbox{}
\vskip 1mm
\noindent
\section*{\lwbt Instructions for Using the Template}
\label{sec:intro}
\vspace{3.5mm}
\noindent
\setlength{\hangindent}{17pt}
1.~You can simply type your text into this template or modify an already prepared manuscript based on the template formatting.

\noindent
\setlength{\hangindent}{17pt}
2.~All Schemes, Equations, Figures, and Tables should be mentioned in the text consecutively and numbered with Arabic numerals, and appear below where they are mentioned for the first time in the main text. To insert Schemes, Equations, Figures, and Tables, create a new break line and paste the resized artworks where desired. Do not modify the amount of space before and after the artworks. One- or two-column format artworks are preferred.


\noindent
\setlength{\hangindent}{17pt}
3.~Save the file with the graphics in place as a Word document.

\noindent
\setlength{\hangindent}{17pt}
4.~Before submitting your manuscript, you need to delete all of the instructions and examples from the template that are not needed. Make sure that all pages have page numbers. Ensure that all parts of the manuscript are presented clearly.

\noindent
\setlength{\hangindent}{17pt}
5.~If it is possible, please provide a set of high resolution Schemes, Equations, Figures and Tables along with your manuscript.


\vskip +10mm
\noindent
\textbf{Manuscript text}. The font is Calibri $\backslash$size\{5\}.~The text may include an introduction, experimental details, theoretical basis, results, discussion, and conclusions.~For articles, please use the decimal system of headings with no more than three levels.
\\


\size{5}
\noindent
Examples:

\vspace{3.5mm}
\section{\lwbt Introduction}
\label{s:introduction}
\noindent
For example:~The parallelization of cutoff pair interactions is mature on CPUs, and typically employs a voxel-based method.

\section{\lwbt Experimental}
\label{s:Experimental}
\noindent
For example: The parallelization of cutoff pair interactions is mature on CPUs, and typically employs a voxel-based method.
\section{\lwbt Results and Discussion}
\noindent
For example: The parallelization of cutoff pair interactions is mature on CPUs, and typically employs a voxel-based method.

\subsection{\lwbt Cutoff pair interactions on GPUs}
\noindent
For example: The parallelization of cutoff pair interactions is mature on CPUs, and typically employs a voxel-based method.

\subsection{\lwbt Cutoff pair interactions on GPUs}
\noindent
For example: The parallelization of cutoff pair interactions is mature on CPUs, and typically employs a voxel-based method.

\subsubsection{\lwbt Cutoff pair interactions on GPUs}
\noindent
For example: The parallelization of cutoff pair interactions is mature on CPUs, and typically employs a voxel-based method.

\subsubsection{\lwbt Cutoff pair interactions on GPUs}
\noindent
For example: The parallelization of cutoff pair interactions is mature on CPUs, and typically employs a voxel-based method.


\begin{lemma} {From $\Omega_{I_{\rm p}}(\mathcal{T},\mathcal{R})$,the probability of computing tag $t_i$'s $K_i$
or the current session key $K_{ji}$ maintained in $r_j$ is negligible.}
\label{lem_history_security}
\end{lemma}

\begin{IEEEproof}
%\textbf{Proof}\quad
The messages in $\Omega_{I_{\rm p}}(\mathcal{T},\mathcal{R})$ that do not involve $K_i$ or $K_{ji}$
do not give any advantage in obtaining these keys.~$K_i$ is used in the form of $E.K_i(K_{ji}||{\rm ID}_{\mathcal{T}}^i)$ and
$K_{ji}$ is used in the form of $E.K_i(K_{ji}||{\rm{ID}}_{\mathcal{T}}^i)$,
$E.K_{ji}(\delta||{\rm ID}_{\mathcal{T}}^i)$, and
$E.K_{ji}(\delta||K'_{ji}||$ $E.K_i$ $(K'_{ji}||
{\textrm{ID}}_{\mathcal{T}}^i))
$.~By Lemma~\ref{lem_history_security}, even if adversaries
obtain many values of these forms, it is computationally
infeasible to obtain $K_i$.~As a result, when the length of $K_i$ and $K_{ji}$ is both $l_{{\rm key}}$,
the probability of computing one of these keys is
$1/2^{l_{{\rm key}}}$.
If $l_{{\rm key}}$ is sufficiently long, this is negligible. Therefore, this lemma holds.
\end{IEEEproof}

\section{\lwbt Conclusions}
\noindent
For example: The parallelization of cutoff pair interactions is mature on CPUs, and typically employs a voxel-based method.


\vskip 6mm
\noindent
References to the literature are cited by \emph{number in square brackets} at appropriate locations (\emph{before} a period, comma, etc.) in the text.

\vskip 6mm
\noindent
Examples:

\setlength{\hangindent}{18pt}
\noindent
  $\bullet$  	Negotiation research spans many disciplines [3].

\setlength{\hangindent}{18pt}
\noindent
  $\bullet$  	This result was later contradicted by Becker and Seligman [5, 6], who .......

\setlength{\hangindent}{18pt}
\noindent
  $\bullet$  	This effect has been widely studied [1-3, 7].

\setlength{\hangindent}{18pt}
\noindent
  $\bullet$  	......achieved until rather recently [11, 21, 22], with......

\setlength{\hangindent}{18pt}
\noindent
  $\bullet$  	......stage of cap formation (see Fig. 5 in Ref. [14]).

\vskip 2mm
\size{5}
\noindent
\textbf{\lwbt Acknowledgment}
\vskip 2mm

\size{5--}
\noindent
The font is Calibri $\backslash$size\{5\--\--\}. Acknowledgements should be inserted at the end of manuscript, before the references. This part should be brief, clear, and should not include any anonyms, titles or effusive comments. Sponsor and financial support acknowledgments are included in the acknowledgment section. For example: This work was supported in part by the US Department of Commerce under Grant BS123456 (sponsor and financial support acknowledgment goes here). Researchers that contributed information or assistance to the article should also be acknowledged in this section.

\vskip 2mm
\size{5}
\noindent
\textbf{\size{5}\lwbt References}
\vskip 2mm

\size{5-}
\noindent
The font is Calibri $\backslash$size\{5\--\}. This part is placed at the end of the manuscript. References should be numbered sequentially as they appear throughout the text. Only one publication should be given for each number. The list of references should only include papers that have been published or accepted by a named publication or recognized preprint server. Authors should ensure the accuracy and completeness of all references before submission. Please ensure references are given in the correct style as shown below in order to avoid delays in typesetting your article
\renewcommand\refname{\size{5}\textbf{\lwbt References}}


\begin{thebibliography}{99}
\size{5-} \addtolength{\itemsep}{-1em}
\vspace {1.5mm}

\bibitem[1]{1}
M. Abrahams and M. Kattenfeld, The role of turbidity as a constraint on predator-prey interactions in aquatic environments, \emph{Behav. Ecol. Sociobiol.,} vol. 40, no. 3, pp.169-174, Mar. 1997. \textbf{(Journal style)}

\bibitem[2]{2}
R. C. Calfee and R. R. Valencia, \emph{APA Guide to Preparing Manuscripts for Journal Publication}. Washington, DC: American Psychological Association, 1991. \textbf{(Authored book style)}

\bibitem[3]{3}
J. M. O'Neill and J. Egan, Men's and women's gender role journeys: metaphor for healing, transition, and transformation, in \emph{Gender Issues Across the Life Cycle}, B. R. Wainrib, Ed. New York: Springer, 1992, pp. 107-123. \textbf{(Book chapter style)}

\bibitem[4]{4}
J. C. Phelan, B. G. Link, A. Stueve, and B. A. Pescosolido, Have public conceptions of mental health changed in the past half century? Does it matter? presented at the 124th Annu. Meeting American Public Health Association, New York, USA, 1969. \textbf{(Paper presented at conferences style)}

\bibitem[5]{5}
A. Mahdavi and B. Spasojevic, Incorporating simulation into building systems control logic, in \emph{Proc. 10th Int. Building Performance Simulation Association Conf.}, Beijing, China, 2007, pp. 1175-1181. \textbf{(Paper in proceedings style)}

\bibitem[6]{6}
X. Yang, Study of building material emissions and indoor air quality, Ph.D. dissertation, Dept. Arch., MIT, MA, USA, 1999. \textbf{(Dissertation style)}

\bibitem[7]{7}
J. Dong, S. Martin, and P. Waldo, Method and system for dynamically presenting cluster analysis results, US Patent US6380937B1, March 30, 2002. \textbf{(Patent style)}

\bibitem[8]{8}
Y. Jiang, Liquid desiccant air-conditioning system and its applications, (in Chinese), \emph{Heating Ventilating \& Air Conditioning}, vol. 34, no. 11, pp. 88-97, 2004. \textbf{(Non-English publication style)}

\bibitem[9]{9}
M. K. Slifka and J. L. Whitton, Clinical implications of dysregulated cytokine production, \emph{J. Mol. Med.}, doi:10.1007/s001090000086. \textbf{(Article by DOI style)}

\bibitem[10]{10}
J. Doe. The dictionary of substances and their effects,  http://www.rsc.org/dose/title, 1999. \textbf{(Online document style)}
  \end{thebibliography}
\begin{strip}
\end{strip}
\clearpage
\size{5}
\noindent
\textbf{FIGURES. }Below each figure, there should be a figure legend. The font is Calibri $\backslash$size\{5-\}. A figure legend normally begins with a brief title describing its whole contents, and continues with a short description of each panel. The format of figure legends is ``Fig. 1\quad Figure legend". The preferred figure formats is TIFF or JPEG with a preferred resolution of \textbf{600dpi} relative to the final figure size. When figures are divided into parts, each part should be labeled with a lower-case (a), (b), and so on, in the same font size as used elsewhere in the figure. All lettering in figures should be in lower-case, except for the first letter of each label which should be capitalized. Use \textbf{a single space} to separate a number and its units. Ensure that the labels are sufficiently large and clear to be readable when the figure is reproduced in the print version of the journal. Figures are referred to in the manuscript as ``Fig. 1" and ``Fig. 1a", or ``Figure 1" at the beginning of a sen-tence. After a manuscript is accepted, we may ask the authors to provide high resolution figures. Failure to supply the file can significantly delay the publication of your work. If possible, please provide a set of high resolution figures along with your initial manuscript. The width of figures should better be 8, 12, or 17 cm, and the height should be less than \mbox{19 cm.}\\

\noindent
\textbf{TABLES.} The font is Calibri $\backslash$size\{5-\}. Above each table, there should be a brief title describing its contents.~Title format should follow~``\textbf{Table 1\quad Table title}".~All details, such as nonstandard abbreviations, and a description of standards of error analysis, should be included in footnotes. Tables are referred to in the manuscript as ``Table 1". The concise type tables are recommended.\\

\noindent
\textbf{EQUATIONS.} The font is Calibri $\backslash$size\{5\}. Equations and mathematical expressions are identified by parenthetical numbers, such as (1), and are referred to in the manuscript as ``Eq. (1)", or ``Equation (1)" at the beginning of a sentence.
%\end{CJK*}
\end{document}


\begin{biography}[yourphotofilename.jpg]
\noindent
\textbf{First A. Author}\ \  Photo. Biographies should be limited to one paragraph consisting of the following: sequentially ordered list of degrees, including years achieved; sequentially ordered places of employ concluding with current employment; associa-tion with any official journals or conferences; major profes-sional and/or academic achievements, i.e., best paper awards, research grants, etc.; any publication information (number of papers and titles of books published); current research interests; association with any professional associations.
\end{biography}

\begin{biography}[yourphotofilename.jpg]
\noindent
\textbf{Second B. Author} Photo. Biographies should be limited to one paragraph consisting of the following: sequentially ordered list of degrees, including years achieved; sequentially ordered places of employ concluding with current employment; associa-tion with any official journals or conferences; major profes-sional and/or academic achievements, i.e., best paper awards, research grants, etc.; any publication information (number of papers and titles of books published); current research interests; association with any professional associations.
\end{biography}
\vskip 22mm
\begin{biography}[yourphotofilename.jpg]
\noindent
\textbf{Third C. Author}  Photo. Biographies should be limited to one paragraph consisting of the following: sequentially ordered list of degrees, including years achieved; sequentially ordered places of employ concluding with current employment; associa-tion with any official journals or conferences; major profes-sional and/or academic achievements, i.e., best paper awards, research grants, etc.; any publication information (number of papers and titles of books published); current research interests; association with any professional associations.
\end{biography} 